\documentstyle[%


righttag%


,11pt%


,amssymb%


,russian%               remove in Eng.


,izvru%                 Set izven% in Eng.


]{amsart}





% Koethe


\newcommand{\tts}[1]{}


\renewcommand{\baselinestretch}{1.03}





\theoremstyle{plain}


\theorembodyfont{\it}


\newtheorem{propnonum}{\hskip\parindent Предложение}


\renewcommand{\thepropnonum}{}








%\hoffset=-15mm%                  For Eng.


\setcounter{page}{29}


\god{2001}


\nomer{10~(473)} %                        Remove in Eng.


%\nomer{Vol.\,45, No.\,10}%               Set in Eng.


%\str{\pageref{begin}--\pageref{end}}%  Set in Eng.


\udk{517.518}





\author{\it М.М.\,Драгилев}


% NB:   ^^ remove in Eng.


\title{О пространствах К\"ете\\


с предэквивалентными безусловными базисами}





\date{}


%\pagestyle{myheadings}%         Set in Eng.


\pagestyle{plain} %              Remove in Eng.


\begin{document}


%\thispagestyle{plain}%          Set in Eng.


\maketitle


\label{begin}





Два базиса в топологическом векторном пространстве называют {\it


предэквивалентными}, если они приводятся к совпадению линейным


автоморфизмом пространства и нормировкой элементов.


Известно, что банахово пространство,


имеющее безусловный базис, лишь в исключительных случаях обладает


тем свойством, что все его безусловные базисы предэквивалентны


\cite{1LT}. В статье доказывается аналогичное утверждение для


метризуемых пространств К\"eте и сопряженных к ним.


\medskip





{\bf 1.} Пусть $\ell$ --- банахово пространство


последовательностей чисел $(t_n)_1^\infty\in {\Bbb C}$ с


безусловным базисом ортов $(e_n)_1^\infty$,


$e_n=(\delta_{kn})_{k=1}^\infty$, с нормой $\|\cdot\|$, монотонной


относительно базиса, удовлетворяющей условию $\|e_n\|=1$


(сокращенно $\ell\in\Lambda$). Матрицу чисел


$$


(a_{pn})_{{\cal P},N}=\{a_{pn}\geq 0: p\in{\cal P},\ n\in \Bbb N\}


$$


называют {\it матрицей К\"eте}, если выполняются следующие условия:


$$


\forall n \mid \sup\limits_p a_{pn}>0;\ \;


\forall p_1,p_2\


\exists p,C>0\mid \max (a_{p_1n},a_{p_2n})\leq Ca_{pn} \ \;


(n=1,2,\dotsc).


$$


Две матрицы называют {\it эквивалентными} ($(a_{pn})_{{\cal


P},N}\sim (b_{qn})_{{\cal Q},N}$), если


$$


\text{а)\ } \forall p\  \exists q,C\mid a_{pn}\leq


Cb_{qn};\quad


\text{б)\ } \forall q\  \exists p,C\mid b_{qn} \leq Ca_{pn}\ \;


(n=1,2,\dotsc).


$$





{\it Пространством К\"eте $($типа $(\ell))$} называют


векторное пространство


$$


\ell((a_{pn})_{{\cal P},N})=\{(t_n)_1^\infty:


(t_na_{pn})_{n=1}^\infty\in \ell,\


\forall\, p\in{\cal P}\}


$$


с локально выпуклой топологией, задаваемой системой преднорм


$$


|(t_n)_1^\infty|_p=\|(t_na_{pn})_{n=1}^\infty\|\quad


(p\in{\cal P}),


$$


где $\ell\in\Lambda$, $\|\cdot\|$ означает норму в $\ell$, а


$(a_{pn})_{{\cal P},N}$ --- матрицу К\"eте.





Эквивалентные матрицы, очевидно, порождают одно и то же


пространство.





В частности, $\ell\in\Lambda$ есть


пространство Кёте:


$\ell=\ell((a_{pn})_{{\cal P},N})$, где $a_{pn}\equiv 1$. Вообще всякое


бесконечномерное банахово пространство $X$, имеющее безусловный


базис $(x_n)_1^\infty$, изоморфно некоторому пространству К\"eте


--- именно, векторному пространству


$$


\ell=\biggl\{(t_n)_1^\infty:\sum_1^\infty t_n\frac{x_n}{\|x_n\|}


\overset{\text{def}}{=} x\in X\biggr\}


$$


с нормой элемента $\|(t_n)_1^\infty\|=\|x\|$. Пространство К\"eте


метризуемо тогда и только тогда, когда его можно задать с помощью


счетной матрицы. Приведем несколько примеров.





Для произвольного множества $A$ числовых последовательностей


обозначим: $|A|=\{(|t_n|)_1^\infty:(t_n)_1^\infty\in


A\}$. Если $\sup\limits_A|t_n|>0$ при всех $n$ и выполнено


условие


$$


\forall\, (t'_n)_1^\infty,(t''_n)_1^\infty\in A\ 


\exists\, (t_n)_1^\infty\in A, C\mid


\max(|t'_n|,|t''_n|)\leq C|t_n|,


$$


то будем рассматривать $|A|$ как матрицу К\"eте.





1) Если $\omega$ и $\varphi$ --- множества (векторные


пространства) всех числовых последовательностей и соответственно


всех финитных числовых последовательностей, то можно считать


$$


\omega=\ell_1(|\varphi|),\quad \varphi=\ell_1(|\omega|).


$$


Пространства К\"eте $\omega$ и $\varphi$ (а также матрицы


$|\omega|$ и $|\varphi|$) называют {\it вырожденными}. Пространство


$\omega$ метризуемо.





2) Пусть $X$ --- отделимое локально выпуклое пространство (л.\,в.\,п.)


с определяющей системой преднорм $(|\cdot|_p)_{\cal P}$.


Последовательность $(x_n)_1^\infty$ элементов $X$, отличных от


нуля, порождает пространство К\"eте (не зависящее от выбора


определяющей системы преднорм в $X$)


$$


\ell ((x_n)_1^\infty)=\ell((|x_n|_p)_{{\cal P},N})\quad


(\ell\in\Lambda).


$$





3) Линейный непрерывный функционал $x'$  в пространстве К\"eте $X$


реализуется в виде числовой последовательности $(u_n)_1^\infty$,


обладающей свойством: ряд вида


$$


\langle x,x'\rangle=\sum_1^\infty t_nu_n \quad (x=(t_n)_1^\infty\in X)


$$


сходится при любом $x$. Как известно \cite{2RR}, система преднорм


$$


|x'|_x =|\langle x,x'\rangle|=\Big|\sum_1^\infty t_nu_n\Big| \quad (x\in X)


$$


задает {\it слабую} топологию $\sigma(X',X)$ в сопряженном к $X$


пространствe $X'$. Положим для произвольного элемента


$x'=(u_n)_1^\infty\in X'$


$$


|x'|_x^*=\sum_1^\infty |t_n|\,|u_n| \quad (x=(t_n)_1^\infty\in X).


$$


Соответствующая топология $\sigma^*(X',X)$ (согласующаяся с той же


двойственностью) мажорирует слабую топологию


$\sigma(X',X)$.





\begin{note}


Если $\ell\in \Lambda$, то топология


$\sigma^*(\ell',\ell)$ в $\ell'$ не сильнее нормированной.


\end{note}





Пространство $X'$ с топологией $\sigma^*(X',X)$ является


пространством К\"eте типа ($\ell_1$):


$$


(X',\sigma^*(X',X))=\ell_1(|X|).


$$





4) Бочечное л.\,в.\,п. $X$ с абсолютным шаудеровским базисом


$(x_n)_1^\infty$ изоморфно пространству К\"eте


$\ell_1((x_n)_1^\infty)$ \cite{3Mit}.





Непосредственно проверяется





\begin{propnonum}


Пространство К\"eте есть отделимое


счетно-полное л.\,в.\,п. с безусловным шаудеровским базисом


$(e_n)_1^\infty$.


\end{propnonum}





В частности, метризуемое пространство К\"eте является


$F$-пространством. Заметим, что не всякое $F$-пространство с


безусловным базисом изоморфно некоторому пространству К\"eте.


\medskip





{\bf 2.} Матрицу К\"eте $(a_{pn})_{{\cal P},N}$ называют





а) {\it нормированной}, если


$$


\exists\, p\mid \inf\limits_n a_{pn}>0 \quad \text{и}\quad


\forall\,p\mid\sup\limits_n a_{pn}<\infty;


$$





б) {\it почти нормированной}, если матрица $(\lambda_n


a_{pn})_{{\cal P},N}$ нормирована при некоторых $\lambda_n>0$;





в) {\it симметричной}, если $(a_{pn})_{{\cal P},N}\sim


(a_{p\sigma(n)})_{{\cal P},N}$ для любой перестановки индексов


$\sigma: \Bbb N\to\Bbb N$.





г) {\it почти симметричной}, если


$$


\forall\,\sigma\ \exists\, (\lambda_n)_1^\infty\mid (a_{pn})_{{\cal P},N}\sim


(\lambda_n a_{p\sigma(n)})_{{\cal P},N}.


$$





\begin{note}


Свойства а), б), в), г) распространяются


на эквивалентные матрицы, а свойства б) и г) --- также на


предэквивалентные матрицы (две матрицы $(a_{pn})_{{\cal P},N}$ и


$(b_{qn})_{{\cal Q},N}$ называют {\it предэквивалентными}, если


$(a_{pn})_{{\cal P},N}\sim (\lambda_n b_{qn})_{{\cal Q},N}$ при некоторых


$\lambda_n$).


\end{note}





\medskip


Для произвольной матрицы К\"eте $(a_{pn})_{{\cal P},N}$ обозначим


\begin{align*}


\delta((a_{pn})_{{\cal P},N})&=\{(t_n)_1^\infty:


\exists\, p_0\ \forall\,p\ \exists\,C\


\forall\, n\,\big|\,|t_n|a_{pn}\leq Ca_{p_0n}\};\\


\delta'((a_{pn})_{{\cal P},N})&=\{(t_n)_1^\infty:


\forall\, p\ \exists\,p_0,\,C\


\forall\, n\,\big|\,|t_n|a_{pn}\leq Ca_{p_0n}\}.


\end{align*}


Например,


$\delta(|\omega|)=\delta(|\varphi|)=\varphi$,


$\delta'(|\omega|)=\delta'(|\varphi|)=\omega$.


Множества $\delta(\cdot)$ и  $\delta'(\cdot)$, очевидно,


сохраняются при переходе к эквивалентным и предэквивалентным


матрицам.





\begin{lem}[{\series{m}\shape{n}\selectfont \cite{4Dr83}}]


$\delta((a_{pn})_{{\cal P},N})\subset \ell_\infty\subset


\delta'((a_{pn})_{{\cal P},N})$.


\end{lem}





{\bf Доказательство.}


Eсли $(t_n)_1^\infty\in


\delta((a_{pn})_{{\cal P},N})$, то


$$


\exists\, p_0\ \forall\,p,\, \exists\,C\


\forall\, n\,\big|\,|t_n|a_{pn}\leq Ca_{p_0n}.


$$


Если $a_{p_0n}\ne 0$, то, полагая $p=p_0$, получим $|t_n|\leq C$.


В противном случае найдется $p$, при котором $a_{pn}\ne 0$. Тогда


будет $t_n=0$, следовательно, $(t_n)_1^\infty\in\ell_\infty$, чем


доказано первое включение. Второе включение вытекает из свойства


матрицы К\"eте


$$


\forall\,p\, \exists\, p_0,C\


\forall\, n\mid a_{pn}\leq Ca_{p_0n}.


$$





\begin{lem}[{\series{m}\shape{n}\selectfont \cite{4Dr83}}]


Матрица $(a_{pn})_{{\cal P},N}$ тогда и


только тогда почти нормирована, когда\linebreak $\delta((a_{pn})_{{\cal


P},N})=\ell_\infty$. При этом пространство К\"eте


$\ell((a_{pn})_{{\cal P},N})$ изоморфно банаховому пространству~$\ell$.


\end{lem}





{\bf Доказательство.} Если при некоторых $\lambda_n$ матрица


$(\lambda_n a_{pn})_{{\cal P},N}$ нормирована, то выполняются


условия


$$


\forall\,p\ \exists\, C\


\forall\, n\mid\lambda_na_{pn}\leq C,\quad


\exists\,p_0, C_0


\mid\inf\limits_n\lambda_na_{p_0n}=C_0>0.


$$


Для произвольной последовательности


$(t_n)_1^\infty\in\ell_\infty$ имеем


$t_n\lambda_na_{pn}\leq \frac{C}{C_0}\lambda_n a_{p_0n}$ или, что


то же самое, $t_na_{pn}\leq \frac{C}{C_0}a_{p_0n}$ ($n=1,2,\dotsc$). Это


означает $(t_n)_1^\infty\in\delta((a_{pn})_{{\cal P},N})$, и


в силу леммы~1\;


$\delta((a_{pn})_{{\cal P},N})=\ell_\infty$. Обратно, если это равенство


выполняется, то, в частности, $(1,1,\dots)\in


\delta((a_{pn})_{{\cal P},N})$. Следовательно,


$\exists\,p_0\ \forall\,p,\ \exists\, C\


\forall\, n\mid a_{pn}\leq Ca_{p_0n}$.


Отсюда можно заключить, что $a_{p_0n}\ne 0$ ($n=1,2,\dotsc$).


Положив $\lambda_n=a_{p_0n}^{-1}$, получаем $\lambda_n


a_{p_0n}=1$, $\lambda_na_{pn}\leq C$ ($n=1,2,\dotsc$), т.\,е.


матрица $(\lambda_na_{pn})_{{\cal P},N}$ нормирована. Так как она


эквивалентна матрице, состоящей из одних единиц, то


$\ell((\lambda_na_{pn})_{{\cal P},N})=\ell$ и, тем более,


пространства $\ell((a_{pn})_{{\cal P},N})$ и $\ell$ изоморфны.





\begin{lem}


Если матрица $(a_{pn})_{{\cal P},N}$ невырождена и


почти симметрична, то $\delta'((a_{pn})_{{\cal


P},N})=\ell_\infty$. При этом либо


$\delta((a_{pn})_{{\cal P},N})\subset c_0$,


либо $\delta((a_{pn})_{{\cal P},N})=\ell_\infty$.


Последнее выполняется тогда и только тогда, когда пространство


К\"eте $\ell((a_{pn})_{{\cal P},N})$ метризуемо.


\end{lem}





{\bf Доказательство.} Если $\delta'((a_{pn})_{{\cal


P},N})\ne\ell_\infty$, то по лемме~1 найдется последовательность


$(t'_{n_k})_1^\infty\in\delta'((a_{pn_k})_{{\cal P},N})$ такая,


что $t'_{n_k}\to\infty$ (при этом, не уменьшая общности, можно


считать, что индексы $n_k$ возрастают, а дополнительная до


$\Bbb N$ последовательность


$(\ov n_k)_1^\infty$ бесконечна). Возможны два случая.





а) Существует последовательность $(t_{n_k})_1^\infty\notin


\delta'((a_{pn_k})_{{\cal P},N})$. Так как $t'_{n_k}\to\infty$,


то можно выделить такую


подпоследовательность $(t'_{n_{k_j}})_1^\infty$, что


$$


|t'_{n_{k_j}}|\geq |t_{n_j}| \quad (j=1,2,\dots).


$$


Пусть $\sigma$ --- перестановка индексов такая, что


$\sigma(n_j)=n_{k_j}$ ($j=1,2,\dotsc$). С одной стороны, имеем


$$


(t'_{n_{k_j}})_1^\infty\in


\delta'((a_{pn_{k_j}})_{{\cal  P},N})=


\delta'((a_{p\sigma(n_j)})_{{\cal P},N}).


$$


С другой стороны,


$$


(t'_{n_{k_j}})_1^\infty\notin \delta'((a_{pn_j})_{{\cal P},N}),


$$


и, таким образом, $\delta'((a_{pn})_{{\cal P},N})\ne


\delta'((a_{p\sigma(n)})_{{\cal P},N})$. Так как это противоречит


предположению о том, что матрица $(a_{pn})_{{\cal P},N}$ почти


симметрична, то случай~а) не имеет места.





б) $(t_{n_k})_1^\infty\in \delta'((a_{pn_k})_{{\cal P},N})$,


какова бы ни была последовательность $(t_{n_k})_1^\infty$, или,


что то же самое, $\delta'((a_{pn_k})_{{\cal P},N})=\omega$. Положив


$\sigma(n_k)=\ov n_k$, $\sigma(\ov n_k)=n_k$ ($k=1,2,\dots$), получим


$$


\delta'((a_{p\ov n_k})_{{\cal P},N})=


\delta'((a_{p\sigma(n_k)})_{{\cal P},N})=


\delta'((a_{pn_k})_{{\cal P},N})=\omega.


$$


Отсюда следует, что $\delta'((a_{pn})_{{\cal P},N})=\omega$. Но


это для почти симметричной матрицы возможно лишь при одном из двух


условий: $(a_{pn})_{{\cal P},N}\sim\omega$ или


$(a_{pn})_{{\cal P},N}\sim |\varphi|$, т.\,е. матрица


является вырожденной. Тем самым доказано, что


$\delta'((a_{pn})_{{\cal P},N})=\ell_\infty$.





Если $\delta((a_{pn})_{{\cal P},N})\not\subset c_0$, то найдется


последовательность $(t_n)_1^\infty\in


\delta((a_{pn})_{{\cal P},N})$, не принадлежащая $c_0$, т.\,е. такая,


что


$$


\underset{n\to\infty}{\varlimsup}|t_n|>0.


$$


Тогда $|t_{n_k}|>\varepsilon$ для некоторого $\varepsilon>0$ и


бесконечной последовательности $(n_k)^\infty_1$ (с бесконечной же


дополнительной последовательностью


$(\ov n_k)^\infty_1$). Пусть $\sigma$ --- перестановка,


меняющая местами $n_k$ и


$\ov n_k$, и пусть $t_n^0=\varepsilon$ ($n=1,2,\dotsc$).


Тогда $(t^0_{n_k})^\infty_1\in\delta((a_{pn_k})_{{\cal P},N})$,


$(t^0_{\ov n_k})_1^\infty=(t^0_{\sigma(n_k)})^\infty_1\in


\delta((a_{p\sigma(n_k)})_{{\cal P},N})=\delta((a_{p\ov n_k})_{\cal P,N})$


и, следовательно, $(t_n)^\infty_1\in\delta((a_{pn})_{{\cal


P},N})$. Отсюда вытекает, что


$\ell_\infty\subset\delta((a_{pn})_{\cal{P},N})$ и в силу леммы~1\;


$\delta((a_{pn})_{{\cal P},N})=\ell_\infty$.





Если пространство К\"eте $\ell((a_{pn})_{{\cal P},N})$ метризуемо,


то можно считать, что числа $\alpha_{pn}$ ($p,n=1,2,\dotsc$) не


убывают по $p$. Так как матрица $(a_{pn})^\infty_{p,n=1}$


невырождена, то для некоторого $p$ найдется последовательность


$(n_k)^\infty_1$ с бесконечным дополнением $(\ov


 n_k)^\infty_1$ и такая, что $a_{pn_k>0}$. Пусть, как и выше,


$\sigma$ --- перестановка, меняющая местами $n_k$ и $\ov n_k$, и


пусть $\lambda_n$ таковы, что матрицы $(a_{pn})^\infty_{p,n=1}$ и


$(\lambda_n a_{pn})^\infty_{p,n=1}$ эквивалентны. Найдутся


$p_1\geqslant p$ и $C$ такие, что $a_{pn}\leqslant C


a_{p_1\sigma(n)}$, и, следовательно, $a_{p_1\ov n_k}\geqslant


C^{-1}a_{pn_k}>0$. Так как и $a_{p_1 n_k}\geqslant a_{pn_k}>0$, то


$a_{p_1n}\ne 0$ при всех $n$. Предположим, что для любого


$p_2>p_1$


$$


\underset{n\to\infty}{\varlimsup}\frac{a_{p_2n}}{a_{p_1n}}<\infty.


$$


Тогда $\delta((a_{pn})^\infty_{p,n=1})=\ell_\infty$. В противном


случае найдутся $p_2$ и подпоследовательность $(n_k)_1^\infty$


такие, что


$$


\lim_{k\to\infty}\frac{a_{p_2n_k}}{a_{p_1n_k}}=\infty.


$$


Имеются две возможности.





а) Для любого $p_3>p_2$


$$


\underset{k\to\infty}{\varlimsup}\frac{a_{p_3n_k}}{a_{p_2n_k}}<\infty.


$$


Тогда $\delta((a_{pn_k})^\infty_{p,k=1})=\ell_\infty$, а поскольку


матрица


$(a_{pn})^\infty_{p,n=1}$ почти симметрична, также и\linebreak


$\delta ((a_{p\ov n_k})^\infty_{p,k=1})=\ell_\infty$,


где $(\ov n_k)^\infty_1$


--- последовательность, дополнительная к $(n_k)^\infty_1$. Следовательно,


$\delta((a_{pn})^\infty_{p,n=1})=\ell_\infty$.





б) Найдутся $p_3$ и последовательность $(n_{k_j})^\infty_1$


такие, что


$$


\lim_{j\to\infty}\frac{a_{p_3n_{k_j}}}{a_{p_2n_{k_j}}}=\infty.


$$


Продолжая этот процесс, либо убедимся после конечного числа шагов


в том, что


$\delta((a_{pn})^\infty_{p,n=1})=\ell_\infty$, либо выделим последовательность


индексов $\nu\subset \Bbb N$ и последовательность


$(p_k)^\infty_1$, обладающие свойством


$$


\lim_{n\in\nu}\frac{a_{p_{k+1},n}}{a_{p_k,n}}=\infty \quad (k=1,2,\dotsc).


$$


Обозначим


$$


t_n=\min \bigg(\frac{a_{p_2n}}{a_{p_1n}}, \frac{a_{p_3n}}{a_{p_2n}}, \dotsc,


\frac{a_{p_{k+1}n}}{a_{p_kn}}\bigg)\ \;


(n\in\nu, \ n_{k-1}\leqslant n<n_k,\ n_0=1, \ k=1,2,\dotsc),


$$


где $n_k\in\nu$ выбираются так, что $t_n>k$ при $n_{k-1}\leqslant


n<n_k$. Если теперь $p$ произвольно и $p_k>p$, то


$$


|t_n|\leqslant \frac{a_{p_{k+1}n}}{a_{p_kn}}\leqslant


\frac{a_{p_{k+1}n}}{a_{pn}},


$$


и, следовательно,


$(t_n)^\infty_1\in\delta'((a_{pn})^\infty_{p,n=1})$. Но это


противоречит доказанному.





Таким образом, для метризуемого пространства К\"eте


$\ell((a_{pn})^\infty_{p,n=1})$ имеем


$\delta((a_{pn})^\infty_{p,n=1})=\ell_\infty$. Обратно, из этого


условия в силу леммы 2 следует, что пространство банахово и, тем


более, метризуемо.


\medskip





{\bf 3.} Как известно \cite{5Dr80}, в вырожденных пространствах


К\"eте $\omega$ и $\varphi$ (и только в них) все безусловные базисы


эквивалентны. Будем рассматривать невырожденные метризуемые


пространства и сопряженные к ним. В основе дальнейшего лежит





\begin{lem}[{\series{m}\shape{n}\selectfont


\cite{1LT}\ (см. также \cite{8Pel}--\cite{12Lin})}]


Банахово


бесконечномерное пространство с предэквивалетными безусловными


базисами изоморфно одному из пространств


$\ell_1$, $\ell_2$, $c_0$.


\end{lem}





\begin{thm}


В метризуемом пространстве К\"eте тогда и только тогда


все безусловные базисы предэквивалентны, когда оно изоморфно


одному из пространств $\omega$,


$\ell_1$, $\ell_2$, $c_0$.


\end{thm}





{\bf Доказательство.} Достаточно рассмотреть невырожденные


пространства. Из того, что в пространстве


$\ell((a_{pn})^\infty_{p,n=1})$ базис ортов предэквивалентен любой


своей перестановке, следует, что матрица $(a_{pn})^\infty_{p,n=1}$


почти симметрична. По лемме 3 заключаем, что


$\delta((a_{pn})^\infty_{p,n=1})=\ell_\infty$. Тогда


$\ell((a_{pn})^\infty_{p,n=1})$ в силу леммы 2


изоморфно пространству $\ell$, а


по лемме 4 --- одному из пространств


$\ell_1$, $\ell_2$, $c_0$. Обратное утверждение содержится в


\cite{5Dr80} и лемме~4.





\begin{lem}


Если $\ell\in\Lambda$ и $\ell'\in\Lambda$, то безусловный


базис в $\ell$ $($соответственно в $\ell')$ является также


безусловным базисом в слабой топологии


$\sigma(\ell,\ell')$ $($соответственно $\sigma(\ell',\ell))$ и наоборот.


\end{lem}





\medskip


Действительно, разложение по базису произвольного элемента


пространства безусловно сходится в слабой топологии, а


единственность разложения следует из того, что коэффициентные


функционалы базиса слабо непрерывны (напр., \cite{2RR},


с.\,55), т.\,е. базис является также слабым базисом. Обратное


утверждение доказывается в (\cite{13Ed}, с.\,622).





\begin{cor}


Если топология $\tau$ в одном из векторных


пространств $\ell$, $\ell'$ слабее исходной, но мажорирует


соответствующую слабую топологию $\sigma(\ell,\ell')$ или


$\sigma(\ell',\ell)$ и если два базиса этого пространства эквивалентны


в топологии $\tau$, то они также эквивалентны в исходной и


слабой топологиях.


\end{cor}





Доказательство следует из того факта, что базисные ряды с


одинаковыми соответствующими коэффициентами одновременно сходятся


или расходятся во всех трех топологиях.





\begin{thm}


Пусть $X'=(X',\sigma^*(X',X))$ --- пространство К\"eте,


сопряженное к невырожденному метризуемому пространству К\"eте


$X=\ell((a_{pn})^\infty_{p,n=1})$, где $\ell'\in\Lambda$ и


$\ell'\ne c_0$. В пространстве $X'$ все безусловные шаудеровские


базисы предэквивалентны тогда и только тогда, когда оно


изоморфно одному из пространств


$(\ell_1,\sigma^*(\ell_1, c_0))$ или $(\ell_2,\sigma^*(\ell_2,\ell_2))$.


\end{thm}





{\bf Доказательство.} Если в пространстве $X'=\ell_1(|X|)$ все


безусловные шаудеровские базисы предэквивалентны, то матрица


$|X|$, а значит, и матрица $(a_{pn})^\infty_{p,n=1}$ почти


симметричны. Применяя леммы 2 и 3, найдем, что пространство $X$


изоморфно $\ell$. В силу леммы 5 и следствия 1 все


безусловные базисы в пространстве $X'=\ell'$ предэквивалентны. По


лемме 4\; $X'$ изоморфно одному из пространств


$\ell_1$, $\ell_2$, $c_0$. Последнее исключается,


т.\,к. в противном случае получили бы


$\ell'=c_0$. Таким образом, $X'$ изоморфно либо пространству


$(\ell_1,\sigma^*(\ell_1,c_0))=\ell_1(|c_0|)$, либо пространству


$(\ell_2,\sigma^*(\ell_2,\ell_2))=\ell_1(|\ell_2|)$.





С другой стороны, все безусловные шаудеровские базисы в каждом из


этих пространств предэквивалентны. Это следует из леммы 4 с учетом


того, что базисы пространств $\ell_1$ и $\ell_2$ являются базисами


также в соответствующих топологиях $\sigma^*$ и обратно, и


аналогичным свойством обладает отношение предэквивалентности


базисов.





\begin{note}


В пространстве $(c_0,\sigma^*(c_0,\ell_1))=\ell_1(|\ell_1|)$,


не являющемся сопряженным к метризуемому пространству К\"eте


\cite{14Kir}, как и в пространствах


$(\ell_1,\sigma^*(\ell_1,c_0))$ и $(\ell_2,\sigma^*(\ell_2,\ell_2))$, все


безусловные шаудеровские базисы предэквивалентны. Доказательство


аналогично предыдущему. То же относится и к каждому из пространств


$(\ell_1,\sigma(\ell_1,c_0))$,


$(\ell_2,\sigma(\ell_2,\ell_2))$, $(c_0,\sigma(c_0,\ell_1))$, сопряженых


соответственно к $\ell_1$, $\ell_2$ и $c_0$, не являющихся


пространствами К\"eте.


\end{note}





\begin{note}


Условие теоремы 2, согласно которому $\ell'\in\Lambda$,


существенно.


\end{note}





{\bf Пример.} В пространстве К\"eте


$(\ell_\infty,\sigma^*(\ell_\infty,\ell_1))$ все безусловные


шаудеровские базисы предэквивалентны.


\medskip





Действительно, шаудеровский безусловный базис в пространстве


$\ell_\infty$ с топологией $\sigma^*(\ell_\infty,\ell_1)$


остается таковым и в топологии $\sigma(\ell_\infty,\ell_1)$.


Обратное также верно, ибо, как нетрудно проверить, сходимость


последовательности элементов пространства $\ell_\infty$ в


топологии $\sigma(\ell_\infty,\ell_1)$ совпадает со сходимостью в


топологии $\sigma^*(\ell_\infty,\ell_1)$. Коэффициентные


функционалы слабого базиса в $\ell_\infty$ образуют слабый базис в


$\ell_1$, а следовательно, и базис в исходной банаховой топологии,


причем данное рассуждение обратимо. Таким образом, существует


взаимно однозначное соответствие между шаудеровскими безусловными


базисами в $(\ell_\infty,\sigma^*(\ell_\infty,\ell_1))$ и


безусловными базисами в $\ell_1$, которое распространяется также


на эквивалентные (предэквивалентные) базисы. Остается применить


лемму 4.





\begin{cor}


В пространстве К\"eте $X'$,


сопряженном к метризуемому пространству К\"eте $X$, все безусловные 


шаудеровские базисы предэквивалентны тогда и только тогда, когда


$X'$ изоморфно одному из пространств $\varphi$,


$(\ell_1,\sigma^*(\ell_1, c_0))$,


$(\ell_2,\sigma^*(\ell_2, \ell_2))$,


$(\ell_\infty,\sigma^*(\ell_\infty, \ell_1))$.


\end{cor}





Целесообразно рассматривать две разновидности пространств К\"eте


$\ell((a_{pn})_{{\cal P},N})$ с почти симметричной матрицей, для


которых соответственно


$\delta((a_{pn})_{{\cal P},N})=\ell_\infty$ и


$\delta((a_{pn})_{{\cal P},N})\ne\ell_\infty$


или, что то же самое, $\delta((a_{pn})_{{\cal


P},N})\subseteq c_0$. Пространства первого рода с предэквивалентными


безусловными базисами исчерпываются банаховыми пространствами


$\ell_1$, $\ell_2$, $c_0$. Пространствами второго рода являются,


в частности, $\omega$ и $\varphi$. Неизвестно, существуют ли


пространства К\"eте второго рода с предэквивалентными безусловными


базисами, алгебраически не изоморфные ни одному из векторных


пространств $\omega$, $\varphi$, $\ell_1$, $\ell_2$, $c_0$,


$\ell_\infty$.





\begin{thebibliography}{99}





\bibitem{1LT}


Lindenstrauss J., Tzafriri L. {\em Classical Banach spaces}. I.


{\em Sequence spaces} // Ergeb. Math. -- 


1977. -- V.\,92. -- 190 p.


\bibitem{2RR}


Робертсон А.П, Робертсон В.Дж. {\em Топологические векторные пространства.}


-- М.: Мир, 1967. -- 257~с.


\bibitem{3Mit}


Митягин Б.С. {\em Аппроксимативная размерность и базисы в ядерных


пространствах} // УМН. -- 1961. -- Т.\,16. -- Вып.\,4. -- С.\,63--132.


\bibitem{4Dr83}


Драгилев М.М. {\em Базисы в пространстве К\"eте}. -- Ростов-на-Дону:


Изд-во ун-та, 1983. -- 144~c.


\bibitem{5Dr80}


Драгилев М.М. {\em Топологические векторные пространства с


эквивалентными базисами} // Матем. заметки. -- 1980. -- Т.\,28. -- Вып.\,6.


-- С.\,945--951.


\bibitem{8Pel}


Pelczinski A. {\em Projection in certain Banach spaces} // Studia Math. --


1960. -- V.\,19. -- P.\,209--228.


\bibitem{9Lorch}


Lorch E. {\em Bicontinuous linear transformation in certain vector


spaces} // Bull. Amer. Math. Soc. -- 1939. -- V.\,45. -- P.\,564--569.


\bibitem{10LinPel}


Lindenstrauss J., Pelczinski A. {\em Absolutely summing operators in


$L_p$-spaces and their applications} // Studia Math. -- 1968. --


V.\,29. -- P.\,275--326.


\bibitem{11Zippin}


Zippin M. {\em On perfectly homogeneus bases in Banach spaces} // Israel


J.  Math. -- 1966. -- V.\,4. -- P.\,265--272.


\bibitem{12Lin}


Lindenstrauss J., Zippin M. {\em Banach spaces with a unique


unconditional basis} // J. Funct. Anal. -- 1969. -- V.\,3. -- P.\,115--125.


\bibitem{13Ed}


Эдвардс Р. {\em Функциональный анализ. Теория и приложения}.


-- М.: Мир, 1969. -- 1071 c.


\bibitem{14Kir}


Кириллов А.А., Гвишиани А.Д. {\em Теоремы и задачи функционального


анализа}. -- М.: Наука, 1989. -- 384 c.





\end{thebibliography}





\vskip 2cm





\post{\it Ростовский государственный}{\it Поступила}


\post{\it университет}{03.07.2000}





\label{end}


\end{document}








\bibitem{6KT}


K\"othe G., Toeplitz O. Lineare R\"aume mit unendlich vielen


Koordinaten und Ringe unendlicher  Mathrizen // J. reine und


angew. Mathem., 1934. V.\,171. S.\,251--270.





\bibitem{7Gel}


Гельфанд И.\,М. Замечание к работе Н.\,К.~Бари <<Биортогональные


системы и базисы в гильбертовом пространстве>> // Уч. зап. Моск.


гос. ун-та, 1945. Т.\,4. В.\,148.





